\subsection{Pretend Polynomials with Rational Exponents} 
We can have expressions that look like polynomials, except that some of the exponents are rational.
\begin{example}
The following expressions have the form of polynomials, except one or more exponents are not positive integers.
\begin{itemize}
\item $x^{2/3}-3x^{1/3}-4$
\item $x^{5/3}-x^{2/3}-20x^{-1/3}$
\item $2x^{-2}-7x^{-1}-15$
\item $4x-8\sqrt{x}+3=\makeblank$
\end{itemize}
\end{example}
Equations with expressions like these can be solved just like polynomial equations (usually by \makeblank[1.5in]), except with an extra final step.
\begin{example}
Solve each equation by treating it as a polynomial equation.
\begin{lpicvsplit}{2}{7}
\foreach \x/\exp in {0/{x^{2/3}-3x^{1/3}-4},1/{x^{5/3}-x^{2/3}-20x^{-1/3}}}
	\regtextcolc{\x}{-1}{$\exp=0$};
\end{lpicvsplit}
\begin{lpicvsplit}{2}{12}
\foreach \x/\exp in {0/{2x^{-2}-7x^{-1}-15},1/{4x-8x^{1/2}+3}}
	\regtextcolc{\x}{-1}{$\exp=0$};
\regtextcoll{1}{-11}{Check:};
\end{lpicvsplit}
\end{example}